\par\addvspace{7pt}\Needspace{4\baselineskip}\noindent\textbf{课堂题：树编码、随机图与随机游走}\par
\begin{exercise}\label{mit:6042-cp25-q2}
\textbf{6.042J，Spring 2015，CP25 第 2 题。}
$n>2$ 个编号顶点的树按如下方式编码：反复删掉\textbf{编号最大的叶}并记录其邻点，剩两个顶点时停止，得到长 $n-2$ 的串。原卷两例的码为 $(6,5,6,2,2)$ 与 $(4,3,2)$。
（a）说明如何由编码重构树；（b）证明它与 $\{1,\ldots,n\}^{n-2}$ 双射，并求编号树数。
\end{exercise}
\sourceof{CP25，物理第 1 页；编者独立解答。原题选最大叶，与本册正文的最小叶版本不同。}
\begin{solution}
(a) 维护尚未删去顶点集 $V$ 和剩余码串 $S$。取 $V$ 中不在 $S$ 出现的最大顶点 $x$，将它与 $S$ 的首项 $y$ 连边，删掉 $x$ 与首项；最后把剩余两点连边。例一依次恢复边 $76,45,56,62,32,12$，例二恢复 $54,43,32,12$。

为什么存在可选点？剩余 $r$ 个顶点时，剩余串长 $r-2$，至多含 $r-2$ 个不同顶点，至少两点未出现，且首项属于 $V$。恢复的边均接一个即将删去的顶点到仍留存的顶点；反向看，它从最后的一条边逐次接叶，所以得到树。

(b) 在编码过程中，每个顶点 $v$ 的出现次数恰为 $d(v)-1$：它作为邻点被记录一次，其度就减一；当自己成为叶被删或成为最后两点之一时，尚余度 $1$。对每个残余树同理，所以“不在剩余串中”恰表示当前叶，最大者正是下一步应删的叶。由此编码树再解码逐步恢复相同的叶边，解码串再编码逐步记录相同首项；两法互逆。共有 $n^{n-2}$ 个码串，因此编号树数为 $n^{n-2}$。这也说明最大叶规则仍得到 Prüfer 双射，不能把正文最小叶的具体解码顺序直接搬来。
\end{solution}

\begin{exercise}\label{mit:6042-cp30-q3}
\textbf{6.042J，Spring 2015，CP30 第 3 题。}
$G$ 为 $n>3$ 顶点简单图，$E(x,y)$ 表示有边，$P(x,y)$ 表示有长 $2$ 的游走。
（a）用 $E$ 写出 $P$ 的谓词公式。以后每对不同顶点独立以概率 $p$ 连边。
（b）对互异 $w,x,y,z$，判定哪些事件对必独立：1 $E(w,x)$ 与 $E(x,y)$；2 $E(w,x)\land E(w,y)$ 与 $E(z,x)\land E(z,y)$；3 $E(x,y)$ 与 $P(x,y)$；4 $P(w,x)$ 与 $P(x,y)$；5 $P(w,x)$ 与 $P(y,z)$。
（c）求 $\Prob(\neg P(x,y))$；（d）令此概率为 $r$，求 $x,y$ 同处某三角形的概率。
\end{exercise}
\sourceof{CP30，物理第 2 页；编者独立解答。}
\begin{solution}
(a) $P(x,y)\iff\exists a\in V(G):E(x,a)\land E(a,y)$；简单图无自环，所以当 $x\ne y$ 时见证点不能是二者之一。
(b) 必独立的是 $1,2,3$。1 使用不同的随机边；2 的四条随机边互异；3 中 $P(x,y)$ 只使用 $xa,ay$（$a\notin\{x,y\}$），不使用边 $xy$，故与其独立。4、5 不保证独立：取 $n=4,p=1/2$，每个长 $2$ 可达事件概率 $7/16$；第 4 对共同发生概率 $17/64$，第 5 对为 $20/64$，均非 $(7/16)^2=49/256$。这些数可逐个数 $2^6$ 个等可能边集：各单事件在 $28$ 个边集中成立，两种交事件分别在 $17,20$ 个边集中成立。端点 $p=0,1$ 的退化情形全部事件当然独立，不改变“必独立”的选择。
(c) 对每个 $a\notin\{x,y\}$，两条边 $xa,ay$ 同在的概率 $p^2$；不同 $a$ 使用不交边组，故
$r=(1-p^2)^{n-2}$。
(d) 同处三角形等价于 $E(x,y)\land P(x,y)$。由 (b) 第 3 项独立，概率为 $p(1-r)$。
\end{solution}

\begin{exercise}\label{mit:6042-cp33-q4}
\textbf{6.042J，Spring 2015，CP33 第 4 题（补充题）。}
$K_n$ 每条边独立染为红、蓝、绿，概率分别 $r,b,g$，和为 $1$。三顶点集合称为三角形，三边同色称为单色。令 $I_T$ 为三角形 $T$ 单色的指示量。
\begin{enumerate}[label=(\alph*)]
\item 求给定三角形单色概率 $m$。
\item 求 $\E I_T$、$\mathrm{Var}(I_T)$。
\item 求两不同三角形同时单色的概率，分别讨论共享一条边和不共享边。
\item 当 $r=b=g=1/3$ 时证明 $I_T,I_U$ 独立。
\item 在此均匀情形，求单色三角形总数 $M$ 的期望与方差。
\item 令 $\mu=\E M$，用 Chebyshev 不等式证明
$\Prob(|M-\mu|>\sqrt{\mu\log\mu})\le1/\log\mu$。
\item 证明上述偏离概率在 $n\to\infty$ 时趋于 $0$。
\end{enumerate}
\end{exercise}
\sourceof{CP33，物理第 2 页；编者独立解答。(f) 补明 $\mu>1$。}
\begin{solution}
(a) 三种单色事件互斥，故 $m=r^3+b^3+g^3$。
(b) $\E I_T=m$，因 $I_T^2=I_T$，方差为 $m-m^2$。
(c) 不共享边时使用不交随机边组，概率 $m^2$；共享一条边时，要同时单色就必须五条不同边全同色，概率 $r^5+b^5+g^5$。
(d) 均匀时 $m=1/9$，共享边的联合概率为 $3(1/3)^5=1/81=m^2$；不共享边已有相同结果。两个 $0$--$1$ 随机变量在取值 $(1,1)$ 的概率等于边际乘积后，另外三种取值概率由减法也等于相应乘积，所以确实独立。
(e) 记 $N=\binom n3$。$\mu=Nm=N/9$，两两独立使所有不同指示量协方差为 $0$，故 $\mathrm{Var}(M)=Nm(1-m)=8N/81=8\mu/9$。这里不需要全部三角形事件相互独立。
(f) 需 $\mu>1$，即整数 $n\ge5$。对非负量 $(M-\mu)^2$，超过 $a^2$ 的概率至多 $\E(M-\mu)^2/a^2$（把期望限制到该事件即可），这就是 Chebyshev 界。因此取 $a=\sqrt{\mu\log\mu}$ 得
\[
\Prob(|M-\mu|>a)\le\frac{8\mu/9}{\mu\log\mu}
=\frac8{9\log\mu}\le\frac1{\log\mu}.
\]
(g) $\mu=\binom n3/9\to\infty$，右端趋于零，概率非负，夹逼即得结论。小 $n$ 不能直接使用含 $\log\mu$ 的表达式。
\end{solution}

\begin{exercise}\label{mit:6042-cp35-q1}
\textbf{6.042J，Spring 2015，CP35 第 1 题。}
随机游走的三幅图用转移矩阵表示，行、列按表中顶点顺序：
\[
P_1=\begin{pmatrix}0&1\\1&0\end{pmatrix}_{x,y},\qquad
P_2=\begin{pmatrix}0&1\\0.9&0.1\end{pmatrix}_{w,z},\qquad
P_3=\begin{pmatrix}1&0&0&0\\1/2&0&1/2&0\\0&1/2&0&1/2\\0&0&0&1\end{pmatrix}_{a,b,c,d}.
\]
\begin{enumerate}[label=(\alph*)]
\item 求图 1 的平稳分布。
\item 解释从 $x$ 出发为何分布不收敛，并刻画哪些初始分布收敛到平稳分布。
\item 求图 2 的平稳分布。
\item 从 $w$ 出发，长期分布是否接近平稳分布？列几步说明。
\item 证明图 3 有不可数多个平稳分布。
\item 从 $b$ 出发，长期在 $d$ 的概率趋近哪个数？
\item 给一个不强连通但只有唯一平稳分布的图。
\end{enumerate}
\end{exercise}
\sourceof{CP35，物理第 1--2 页，图 1--3；编者独立解答。在线 Random Walks 同图，另保留各 Q 编号。}
\begin{solution}
平稳分布指概率行向量 $\pi$ 满足 $\pi P=\pi$。
(a) 方程为 $\pi_x=\pi_y$，故 $(1/2,1/2)$。
(b) 初始 $(q,1-q)$ 每步交换两项。仅当 $q=1/2$ 才收敛（且始终平稳），否则奇偶步有不同子序列极限；从 $x$ 出发就是 $(1,0),(0,1)$ 交替。
(c) 方程 $\pi_w=0.9\pi_z$ 和两项和 $1$，给 $(9/19,10/19)$。
(d) 记第 $t$ 步在 $w$ 的概率为 $q_t$，则 $q_{t+1}=0.9(1-q_t)$，从 $q_0=1$ 得
$q_t=9/19+(10/19)(-0.9)^t\to9/19$。前几项为 $1,0,0.9,0.09,0.819$，振荡幅度逐步衰减，所以会接近平稳分布。
(e) 平稳方程给 $\pi_b=\pi_c/2$、$\pi_c=\pi_b/2$，故二者为 $0$。全部平稳分布是 $(s,0,0,1-s)$，$0\le s\le1$；参数区间不可数。
(f) 令从 $b,c$ 最终吸收到 $d$ 的概率分别为 $h_b,h_c$。第一步分解给 $h_b=h_c/2$、$h_c=(h_b+1)/2$，所以 $h_b=1/3$。每一步仍留在 $\{b,c\}$ 的条件概率为 $1/2$，故 $t$ 步未吸收概率为 $2^{-t}$；因此在 $d$ 的概率确实趋于最终吸收概率 $1/3$。
(g) 两点 $u\to v$ 概率 $1$，$v\to v$ 概率 $1$。只有 $(0,1)$ 平稳，且 $v$ 不可达 $u$，不强连通。
\end{solution}

\begin{exercise}\label{mit:6042-cp35-q2}
\textbf{6.042J，Spring 2015，CP35 第 2 题。}
证明有限非空随机游走图的均匀分布平稳，当且仅当每个顶点的所有入弧概率之和为 $1$。
\end{exercise}
\sourceof{CP35，物理第 2 页；编者独立解答。}
\begin{solution}
设有 $n\ge1$ 个顶点，转移矩阵为 $P=(p_{uv})$，缺弧的概率视为 $0$。均匀行向量 $\pi$ 的每项为 $1/n$。对每个 $v$，$(\pi P)_v=(1/n)\sum_u p_{uv}$。这等于 $\pi_v=1/n$ 当且仅当列和为 $1$。逐个 $v$ 成立即为所需等价。
\end{solution}

\begin{exercise}\label{mit:6042-cp35-q3}
\textbf{6.042J，Spring 2015，CP35 第 3 题。}
Google-graph 指每个顶点的出弧等概率选择的有向随机游走图。图是对称的，即 $v\to w$ 当且仅当 $w\to v$。设每点出度正、总弧数 $e>0$，定义 $d(v)=d^+(v)/e$。
（a）若以 $d$ 作网页排名，怎样人为提高自己页面的分数？为什么这种简单方法不能直接用于一般有向图的实际排名模型？（b）证明 $d$ 平稳。
\end{exercise}
\sourceof{CP35，物理第 2 页；编者独立解答，补明正出度条件；(a) 仅讨论原题的简化模型。}
\begin{solution}
(a) 添加很多专门连接自己页面的辅助页面，并使链接成对出现，会提高自己页面的度数。在由中心页面与 $k$ 个辅助页构成的星形图中，中心分数为 $k/(2k)=1/2$，与原本很小的分数相比可大幅提高。一般有向图不保证反向弧存在，平稳权重取决于链接的方向及来源页的概率流，不能只凭本地度数计算；这解释了原简化攻击为何不能据此推出实际排名分数。这里不将历史玩具模型当作当前搜索系统规则。
(b) 度数和为 $e$，故 $d$ 是概率分布。对每个 $v$，
\[
(dP)(v)=\sum_{u:u\to v}\frac{d^+(u)}e\frac1{d^+(u)}
=\frac{d^-(v)}e=\frac{d^+(v)}e=d(v),
\]
最后等式用对称性。因此平稳。若有零出度顶点，转移概率 $1/d^+(u)$ 未定义，须先另定该点的转移规则，不能套此原式。
\end{solution}
